% Wave Motions in the Ocean: Myrl's View
% Chapter 4; historical source: source/ChapmanRizzoli4.pdf.
% Source-page comments preserve printed/physical page correspondence.
% Substantive corrections are documented in ERRATA.md.

\setcounter{chapter}{3}
\setcounter{page}{64}

% -----------------------------------------------------------------------------
% Source printed page 64 / physical page 1
% -----------------------------------------------------------------------------
\chapter{Internal gravity waves}

We have seen that gravity can provide the restoring force which allows waves to
propagate along an interface. We saw that if the interface separates two fluids with
slightly different densities, then a much slower version of surface waves -- called internal
waves -- is possible. We now turn our attention to a more thorough investigation of
internal gravity waves. In particular, we extend the previous ideas to situations in
which the vertical density stratification varies continuously within the fluid. And, since
internal waves are ubiquitous in the ocean, we will spend some time reviewing their
observed properties, as well.

\section{The internal wave equation}

Suppose we do the following thought experiment. In a continuously stratified fluid, we
raise a parcel of water from its equilibrium position a small amount $\xi$.

\sourcepagebreak

% -----------------------------------------------------------------------------
% Source printed page 65 / physical page 2
% -----------------------------------------------------------------------------
% Vector reconstruction; source PDF ChapmanRizzoli4.pdf, physical page 2, trim=185bp 445bp 155bp 75bp.
\wavevectorart[0.55]{ch04-p065-parcel-displacement}

The change in pressure experienced by the parcel is $dp=-\rho_0g\xi$ while the change in
density is $d\rho=dp/c^2$. At this point, the buoyancy force acting on the parcel (per unit
volume) introduces an acceleration, so that
\[
 g(\rho_{out}-\rho_{in})
 =g[(\rho_0+\rho_{0z}\xi)-(\rho_0-\rho_0g\xi/c^2)]
 =\rho_0\xi_{tt}.
\]
Rearranging
\[
 \xi_{tt}+\xi\left(-\frac{g\rho_{0z}}{\rho_0}-\frac{g^2}{c^2}\right)=0
\]
which is a simple harmonic oscillator equation with solution $e^{\pm iNt}$ where
\[
 N^2(z)=-\frac{g\rho_{0z}}{\rho_0}-\frac{g^2}{c^2}.
\]
Thus the parcel oscillates about its equilibrium position at a natural frequency
determined by the local density stratification and the fluid's compressibility. This
frequency is called the buoyancy, Brunt--V\"ais\"al\"a or V\"ais\"al\"a frequency and is often used
to characterize the degree of stratification in the ocean.

For applications to the ocean, the effect of compressibility is typically neglected
because $g^2/c^2$ is usually small compared to $g\rho_{0z}/\rho$, so we will neglect it here. In the
atmosphere, compressibility is often important, so the full definition of $N^2$ must be
used. A brief discussion of this case is presented by Gill (1982, pp. 169--175).

\sourcepagebreak

% -----------------------------------------------------------------------------
% Source printed page 66 / physical page 3
% -----------------------------------------------------------------------------
The momentum, mass and continuity equations for a rotating, incompressible
fluid are
\[
 \frac{D\vec u^{\,*}}{Dt}+f\hat k\times\vec u^{\,*}
 =-\frac{\nabla p^*}{\rho^*}-g\hat k
\]
\[
 \frac{\partial\rho^*}{\partial t}+\nabla\cdot\rho^*\vec u^{\,*}=0
\]
\[
 \frac{D\rho^*}{Dt}=0\ \Rightarrow\ \nabla\cdot\vec u^{\,*}=0.
\]
One solution to these equations is a motionless, hydrostatic balance, i.e.
$\vec u_0=0$; $0=-p_{0z}-g\rho_0(z)$. Each dynamic variable may then be separated into a
hydrostatic part and a small departure from it
\[
 \vec u^{\,*}=\vec u_0+\vec u\ ;\qquad p^*=p_0+p\ ;\qquad \rho^*=\rho_0+\rho.
\]
After substituting these into the full equations, assuming that the departures are very
small perturbations and neglecting the nonlinear terms, we obtain
\[
 \frac{\partial\vec u}{\partial t}+f\hat k\times\vec u
 =-\frac{\nabla p}{\rho_0}-\frac{g\rho\hat k}{\rho_0}
\]
\[
 \rho_t+w\rho_{0z}=0
\]
\[
 \nabla\cdot\vec u=0.
\]
Next we specialize to periodic motion $e^{-i\sigma t}$ and write out components
\begin{align*}
 -i\sigma u-fv&=-p_x/\rho_0,\\
 -i\sigma v+fu&=-p_y/\rho_0,\\
 -i\sigma w&=-p_z/\rho_0-g\rho/\rho_0,\\
 u_x+v_y+w_z&=0,\\
 -i\sigma\rho+w\rho_{0z}&=0.
\end{align*}
Eliminate $\rho$ between the vertical momentum equation and the density equation
\[
 \rho_0(N^2-\sigma^2)w=i\sigma p_z.
\]

\sourcepagebreak

% -----------------------------------------------------------------------------
% Source printed page 67 / physical page 4
% -----------------------------------------------------------------------------
The horizontal momentum equations may be rewritten
\[
 u=\frac{1}{\rho_0}\frac{-i\sigma p_x+fp_y}{\sigma^2-f^2}
\]
\[
 v=\frac{1}{\rho_0}\frac{-i\sigma p_y-fp_x}{\sigma^2-f^2}
\]
from which continuity becomes
\[
 -i\sigma\nabla_H^2p+(\sigma^2-f^2)\rho_0w_z=0
\]
where $\nabla_H^2=\partial^2/\partial x^2+\partial^2/\partial y^2$. Now eliminate the pressure to obtain
\[
 \rho_0(N^2-\sigma^2)\nabla_H^2w-(\sigma^2-f^2)(\rho_0w_z)_z=0
\]
or
\[
 \frac{1}{\rho_0}(\rho_0w_z)_z
 -\left(\frac{N^2-\sigma^2}{\sigma^2-f^2}\right)\nabla_H^2w=0.
\]
Now if we had let $\rho_0$ be constant in the momentum equations [so that it is
differentiated only in $N^2(z)$], then we would have obtained
\[
 w_{zz}-\frac{N^2-\sigma^2}{\sigma^2-f^2}\nabla_H^2w=0.
\]
This last simplification is called the \emph{Boussinesq approximation} and it means that
$\rho_{0z}/\rho_0\ll w_z/w$, i.e. $\rho_0(z)$ changes over a large vertical scale. It is quite adequate in
the ocean.

At the free surface (very close to $z=0$), we have $Dp^*/Dt=0$. Let $p^*=p_0+p$,
linearize and apply the result at $z=0$:
\[
 -i\sigma p+w p_{0z}=0\qquad\text{at}\qquad z=0
\]
or
\[
 -i\sigma p-wg\rho_0=0\qquad\text{at}\qquad z=0.
\]
Now use the previous equations to eliminate $p$
\[
 (\sigma^2-f^2)w_z+g\nabla_H^2w=0\qquad\text{at}\qquad z=0.
\]

\sourcepagebreak

% -----------------------------------------------------------------------------
% Source printed page 68 / physical page 5
% -----------------------------------------------------------------------------
At a flat bottom
\[
 w=0\qquad\text{at}\qquad z=-D.
\]
So we have to solve

% Vector reconstruction; source PDF ChapmanRizzoli4.pdf, physical page 5, trim=105bp 385bp 90bp 215bp.
\wavevectorart[0.72]{ch04-p068-boundary-value-problem}

\section{Unbounded, rotating, stratified fluid}

We suppose for the moment that $N^2(z)=\text{constant}$. We also assume that the
coordinates rotate around the $z$ axis at $\Omega$ so that $f=2\Omega=\text{constant}$ which is the
$f$-plane approximation. The field equation is
\[
 w_{zz}-\left(\frac{N^2-\sigma^2}{\sigma^2-f^2}\right)(w_{xx}+w_{yy})=0.
\]
Since $N,f$ are constants, exact solutions are
\[
 w=e^{-i\sigma t+ikx+i\ell y+imz}
\]
from which
\[
 m^2=\left(\frac{N^2-\sigma^2}{\sigma^2-f^2}\right)(k^2+\ell^2)
\]
is the dispersion relation. In $k,\ell,m$ space, the dispersion relation is a cone if
$f^2<\sigma^2<N^2$ or $f^2>\sigma^2>N^2$.

\sourcepagebreak

% -----------------------------------------------------------------------------
% Source printed page 69 / physical page 6
% -----------------------------------------------------------------------------
% Vector reconstruction; source PDF ChapmanRizzoli4.pdf, physical page 6, trim=190bp 570bp 160bp 60bp.
\wavevectorart[0.45]{ch04-p069-dispersion-cone}

All possible wave vectors for waves of frequency $\sigma$ lie on this cone. They may have
any length. Fixing $\sigma$ fixes their direction. If we define $\theta$ as in the sketch, then with
$K=(k^2+\ell^2+m^2)^{1/2}$, we can write
\[
 m=K\cos\theta\ ;\qquad (k^2+\ell^2)^{1/2}=K\sin\theta
\]
and the dispersion relation may be rewritten as
\[
 \sigma^2=N^2\sin^2\theta+f^2\cos^2\theta
\]
or
\[
 \sigma^2K^2=N^2(k^2+\ell^2)+f^2m^2.
\]
These waves are of the form
$(\vec u,p,\rho)=(\vec u_0,p_0,\rho_0)e^{-i\sigma t+i\vec k\cdot\vec x}$. By continuity,
$\nabla\cdot\vec u$ leads to $\vec k\cdot\vec u_0=0$ which means that the fluid motion occurs in planes \emph{normal} to the
wave vector. That is, the waves are transverse. Also, $\nabla p\sim\vec k p_0$ which is perpendicular
to $\vec u$, so the pressure gradient forces are normal to fluid flow and acceleration.

If $f=0$, then the momentum equation in any direction normal to $\vec k$

\sourceart[0.40]{ChapmanRizzoli4.pdf}{6}{220bp 60bp 195bp 610bp}

becomes
\[
 u_t=-gp\sin\theta/\rho_0
\]

\sourcepagebreak

% -----------------------------------------------------------------------------
% Source printed page 70 / physical page 7
% -----------------------------------------------------------------------------
while
\[
 \rho_t+u\sin\theta\,\rho_{0z}=0
\]
because the pressure gradient forces are along $\vec k$. This says that only the density
gradient along $\vec u$ matters, from which
\[
 u_{tt}+N^2\sin^2\theta\,u=0\ ;\qquad \sigma^2=N^2\sin^2\theta
\]
and motion is along a straight line perpendicular to $\vec k$.

If $N=0$, then the momentum equation in any direction normal to $\vec k$

\sourceart[0.46]{ChapmanRizzoli4.pdf}{7}{185bp 350bp 170bp 240bp}

becomes
\[
 \vec u_t+f\cos\theta\,\hat k'\times\vec u=0.
\]
Thus the motion occurs in inertial circles at $\sigma^2=f^2\cos^2\theta$.

We can examine the group velocity by defining
\[
 W(k,\ell,m,\sigma)
 =m^2-\frac{N^2-\sigma^2}{\sigma^2-f^2}(k^2+\ell^2)=0.
\]
Now
\[
 dW\big|_{\ell,m}
 =\frac{\partial W}{\partial k}\,dk
 +\frac{\partial W}{\partial\sigma}\,d\sigma=0
\]

\sourcepagebreak

% -----------------------------------------------------------------------------
% Source printed page 71 / physical page 8
% -----------------------------------------------------------------------------
so
\[
 \left.\frac{\partial\sigma}{\partial k}\right|_{\ell,m}
 =-\frac{\left.\partial W/\partial k\right|_{\ell,m}}
         {\left.\partial W/\partial\sigma\right|_{\ell,m}}
 \qquad\text{etc.}
\]
From this
\begin{align*}
 c_{gx}&=\frac{\partial\sigma}{\partial k}
 =\frac{k}{\sigma}\frac{N^2-\sigma^2}{K^2},\\
 c_{gy}&=\frac{\partial\sigma}{\partial\ell}
 =\frac{\ell}{\sigma}\frac{N^2-\sigma^2}{K^2},\\
 c_{gz}&=\frac{\partial\sigma}{\partial m}
 =-\frac{m}{\sigma}\frac{\sigma^2-f^2}{K^2}.
\end{align*}
First notice that
\begin{align*}
 \vec k\cdot\vec c_g
 &=\left(\frac{k^2}{\sigma}+\frac{\ell^2}{\sigma}\right)
 \left(\frac{N^2-\sigma^2}{K^2}\right)
 -\frac{m^2}{\sigma}\left(\frac{\sigma^2-f^2}{K^2}\right)\\
 &=\frac{-\sigma^2(k^2+\ell^2+m^2)+N^2(k^2+\ell^2)+f^2m^2}{K^2\sigma}\\
 &=0.
\end{align*}
This means that the group velocity is perpendicular to the phase velocity! Next
observe that
\[
 \vec c_g=\frac{1}{\sigma K^2}
 [\hat i k(N^2-\sigma^2)+\hat j\ell(N^2-\sigma^2)
 +\hat k m(f^2-\sigma^2+N^2-N^2)]
\]
That is
\[
 \vec c_g=\frac{1}{\sigma K^2}
 [\vec k(N^2-\sigma^2)-\hat k m(N^2-f^2)].
\]
This allows us to visualize the direction of $\vec c_g$.

\sourcepagebreak

% -----------------------------------------------------------------------------
% Source printed page 72 / physical page 9
% -----------------------------------------------------------------------------
\sourceart[0.76]{ChapmanRizzoli4.pdf}{9}{90bp 470bp 80bp 60bp}

Finally, using the dispersion relation,
\begin{align*}
 |\vec c_g|
 &=\left(c_{gx}^2+c_{gy}^2+c_{gz}^2\right)^{1/2}\\
 &=\frac{[k^2(N^2-\sigma^2)^2+\ell^2(N^2-\sigma^2)^2
 +m^2(\sigma^2-f^2)^2]^{1/2}}{\sigma K^2}\\
 &=\frac{[(N^2-\sigma^2)^2\sin^2\theta
 +(\sigma^2-f^2)^2\cos^2\theta]^{1/2}}{\sigma K}\\
 &=\frac{[(N^2-f^2)^2\sin^4\theta\cos^2\theta
 +(N^2-f^2)^2\cos^4\theta\sin^2\theta]^{1/2}}{\sigma K}\\
 &=\frac{|N^2-f^2|\sin\theta\cos\theta}{\sigma K}.
\end{align*}
An alternative is to define $\phi$ as the angle of energy propagation such that
\[
 \sigma^2=N^2\cos^2\phi+f^2\sin^2\phi
\]
\[
 |\vec c_g|=|N^2-f^2|\sin\phi\cos\phi/(\sigma K).
\]

\sourceart[0.70]{ChapmanRizzoli4.pdf}{9}{105bp 95bp 95bp 520bp}

There is really nothing mysterious about $\vec c_g,\vec c_{ph}$. For deep water waves we had

\sourcepagebreak

% -----------------------------------------------------------------------------
% Source printed page 73 / physical page 10
% -----------------------------------------------------------------------------
$c_g=\frac12c_{ph}$ which states that, in a group, individual crests arise at the trailing end,
propagate through the group faster than the group goes, and die out at the leading
edge. For these internal gravity waves, individual crests arise at one side of the group
and move through it at right angles to the group motion, finally dying out at the other
side. An example is depicted by Gill (1982, pp. 135--6).

Imagine a harmonic source $e^{-i\sigma t}$ at the origin of space coordinates and let's
discuss the wave field for various choices of $\sigma,N,f$. In the general case, energy is
localized to the cone whose apex angle is $\phi$ defined by
$\sigma^2=N^2\cos^2\phi+f^2\sin^2\phi$.

(i) Rotation only: $N=0$, $f\ne0$. We have
\[
 \sigma=f\sin\phi\ ;\qquad |\vec c_g|=f\cos\phi/K.
\]

\sourceart[0.64]{ChapmanRizzoli4.pdf}{10}{150bp 305bp 135bp 255bp}

This says that for vertical energy propagation, $\phi=0^\circ$, the flow is steady ($\sigma=0$) but
the group velocity is nonzero. The wavenumbers are arbitrary but horizontal, i.e.
$m=0$. The flow is basically that of Taylor columns which are steady, geostrophic flows
with no vertical variability ($\partial/\partial z=0$). If the direction of energy propagation is
horizontal, $\phi=90^\circ$, then the frequency must be the inertial frequency ($\sigma=f$) and the
group velocity is zero. The wavenumber is purely vertical, i.e. $k=\ell=0$. This
corresponds to solutions of the horizontal momentum equations in which $p$ is a
function of $z$ only. The pressure gradient terms disappear leaving
\[
 u_t-fv=0
\]

\sourcepagebreak

% -----------------------------------------------------------------------------
% Source printed page 74 / physical page 11
% -----------------------------------------------------------------------------
\[
 v_t+fu=0
\]
which has the solution $\sigma=f$, $u=iv$. Thus, we see that in a rotating homogeneous
fluid, low frequency energy flows vertically in the form of Taylor columns, while inertial
oscillations at different depths are entirely independent.

(ii) Stratification only: $f=0$, $N\ne0$. We have
\[
 \sigma=N\cos\phi\ ;\qquad |\vec c_g|=N\sin\phi/K.
\]

\sourceart[0.64]{ChapmanRizzoli4.pdf}{11}{150bp 320bp 130bp 230bp}

This says that for vertical energy propagation, $\phi=0^\circ$, the flow oscillates at the
buoyancy frequency ($\sigma=N$) and the group velocity is zero. The wavenumber is
arbitrary and horizontal, i.e. $m=0$. These are called buoyancy oscillations. The flow
has Taylor column-like structure, columnar in the vertical, but it is like inertial
oscillations in that energy does not propagate. For horizontal energy propagation,
$\phi=90^\circ$, the flow is steady ($\sigma=0$) and the group velocity is nonzero. The wavenumber
is vertical. In this case, the momentum equations reduce to
$0=-\nabla p/\rho_0-\hat k g\rho/\rho_0$ from which $\rho$ must be zero since it does not vary in time
and can be absorbed into $\rho_0$. This leads to $w=0$ from the density equation, leaving
$u_x+v_y=0$ from continuity. So, each layer in the stratified flow moves independently
of all others. Flow in each layer is nondivergent and buoyancy has no effect. In two
dimensions, if $v_y=0$ ($v=0$, say), then $u_x=0$, i.e. $u=u(z)$. Thus, low frequency
energy flows horizontally and, in two dimensions, is analogous to the Taylor column flows.

\sourcepagebreak

% -----------------------------------------------------------------------------
% Source printed page 75 / physical page 12
% -----------------------------------------------------------------------------
Notice quite generally that effects due to $f$ and effects due to constant $N$ are
very similar in their mathematical expression.

(iii) Both rotation and constant stratification:
\[
 \sigma^2=N^2\cos^2\phi+f^2\sin^2\phi
\]
\[
 |\vec c_g|=|N^2-f^2|\sin\phi\cos\phi/(\sigma K).
\]

\sourceart[0.62]{ChapmanRizzoli4.pdf}{12}{145bp 360bp 135bp 180bp}

For vertical energy propagation, we recover the buoyancy oscillations while for
horizontal energy propagation, we recover the inertial oscillations. Now there are no
zero frequency wave flows that propagate energy. At frequency $\sigma$, energy is confined to
the cone whose sides lie at $\phi$ to the vertical.

What happens if the source frequency is outside the range of $f$ to $N$? In that
case the field equation can be written
\[
 w_{zz}+\frac{\sigma^2-N^2}{\sigma^2-f^2}(w_{xx}+w_{yy})=0
\]
and we see that the very nature of the equation has changed from a hyperbolic (or
wave-type) equation to an elliptic (or potential flow type) equation because the sign of
the coefficient has changed. Our free-wave dispersion relation now becomes
\[
 m^2=-\frac{\sigma^2-N^2}{\sigma^2-f^2}(k^2+\ell^2).
\]

\setcounter{page}{76}

% -----------------------------------------------------------------------------
% Source printed page 76 / physical page 13
% -----------------------------------------------------------------------------
and we see that at least one of the wavenumbers must be complex. This, in turn,
means that the solution is no longer free to propagate, but instead must decay
exponentially in some direction. For example, for waves propagating in the horizontal
($k,\ell$ are real), the oscillations must grow or decay monotonically in the vertical.

\section{Waveguide modes}

We turn our attention to the full problem introduced at the beginning of the chapter.
Instead of an unbounded fluid, there are now surface and bottom boundaries to
contend with.

\sourceart[0.72]{ChapmanRizzoli4.pdf}{13}{105bp 330bp 90bp 245bp}

Initially we restrict ourselves to $N^2=\text{constant}$. The signs of
\[
 S^2\equiv\sigma^2-f^2\ ;\qquad
 R^2\equiv\frac{N^2-\sigma^2}{\sigma^2-f^2}
\]
are crucial and we have several cases to consider.

\sourcepagebreak

% -----------------------------------------------------------------------------
% Source printed page 77 / physical page 14
% -----------------------------------------------------------------------------
\sourceart[0.78]{ChapmanRizzoli4.pdf}{14}{85bp 495bp 75bp 45bp}

Evidently the cases $\sigma^2>N^2,f^2$ and $\sigma^2<N^2,f^2$ are identical for either
$N^2>f^2$ or $N^2<f^2$, but the case where $\sigma^2$ is intermediate between $N^2$ and $f^2$ depends strongly
on whether $N^2>f^2$ or $N^2<f^2$. We will look at cases A and B separately.

Case A: Define $R_1^2=(\sigma^2-N^2)/(\sigma^2-f^2)$ and consider $R_1^2>0$. We let
\[
 w=e^{-i\sigma t+ikx}\omega(z)
\]
and $\omega$ satisfies
\[
 (\sigma^2-f^2)\omega_z-gk^2\omega=0\qquad z=0
\]
\[
 \omega_{zz}-k^2R_1^2\omega=0
\]
\[
 \omega=0\qquad z=-D.
\]
Solutions are
\[
 w=e^{-i\sigma t+ikx}\sinh[kR_1(z+D)]
\]
and they satisfy the top ($z=0$) boundary condition only if
\[
 \frac{R_1(\sigma^2-f^2)}{gk}=\tanh(kR_1D).
\]
This last relation is effectively the dispersion relation; its solution is $\sigma=\sigma(k)$ or, as we
have done the problem, $k=k(\sigma)$. To see what solutions exist, plot the left-hand side
and the right-hand side versus $k$.

\sourcepagebreak

% -----------------------------------------------------------------------------
% Source printed page 78 / physical page 15
% -----------------------------------------------------------------------------
% Vector reconstruction; source PDF ChapmanRizzoli4.pdf, physical page 15, trim=190bp 615bp 165bp 60bp.
\wavevectorart[0.45]{ch04-p078-case-a-intersections}

This is case A with $S^2=\sigma^2-f^2>0$. There are two oppositely travelling waves which
we can identify with the usual surface waves existing in the absence of stratification
and rotation.
\[
 N^2=f^2=0\ ;\qquad R_1^2=1\ ;\qquad \sigma^2-f^2=\sigma^2>0.
\]
Notice that in case A1 ($\sigma^2<f^2,N^2$), no waves exist. The plot of the left-hand side
and the right-hand side versus $k$ looks like

% Vector reconstruction; source PDF ChapmanRizzoli4.pdf, physical page 15, trim=190bp 300bp 165bp 390bp.
\wavevectorart[0.45]{ch04-p078-case-a1-no-intersections}

and there are no solutions to the proposed dispersion relation. What has happened is
that our assumption of free wave propagation in the $x$ direction (real $k$) has proved
impossible to satisfy.

Case B: Now $R^2=(N^2-\sigma^2)/(\sigma^2-f^2)>0$. We let
\[
 w=e^{-i\sigma t+ikx}\omega(z)
\]
and $\omega$ satisfies
\[
 (\sigma^2-f^2)\omega_z-gk^2\omega=0\qquad z=0
\]
\[
 \omega_{zz}+k^2R^2\omega=0.
\]

\sourcepagebreak

% -----------------------------------------------------------------------------
% Source printed page 79 / physical page 16
% -----------------------------------------------------------------------------
\[
 \omega=0\qquad z=-D.
\]
Solutions this time are
\[
 w=e^{-i\sigma t+ikx}\sin[kR(z+D)]
\]
with
\[
 \frac{R(\sigma^2-f^2)}{gk}=\tan(kRD).
\]
Again we look at the dispersion relation

\sourceart[0.82]{ChapmanRizzoli4.pdf}{16}{70bp 420bp 55bp 125bp}

In both cases there is now an infinite set of oppositely travelling modes $n=1,2,\ldots$. The
case $\sigma^2>f^2$ has an additional pair of small-$k$ modes not present in the case $\sigma^2<f^2$.

For large $k$, the $n=1,2,\ldots$ modes have the approximate dispersion relation
\[
 k_nRD=\pm n\pi,
\]
or
\[
 k_nD\left(\frac{N^2-\sigma^2}{\sigma^2-f^2}\right)^{1/2}=\pm n\pi.
\]
This does not hold for the small $k$ ($n=0$) modes. For them, if $kRD\ll1$, then we
obtain
\[
 \frac{\sigma^2-f^2}{gD}=k_0^2.
\]
These we recognize as the old surface modes in shallow water now modified by
rotation. Note that they do not exist when $\sigma^2<f^2$.

\sourcepagebreak

% -----------------------------------------------------------------------------
% Source printed page 80 / physical page 17
% -----------------------------------------------------------------------------
Notice that if we require $w=\omega=0$ at $z=0$, i.e. a rigid lid, then the dispersion
relation is $\sin(kRD)=0$, so that $k_nRD=\pm n\pi$ becomes exact. But we no longer have
the surface modes $k_0$.

Let's look at the dispersion relations more closely. For the surface modes
\[
 \sigma^2\simeq k^2gD+f^2.
\]
For the internal modes
\[
 (\sigma^2-f^2)(n\pi/kD)^2\simeq(N^2-\sigma^2)
\]
\[
 \sigma^2[1+(n\pi/kD)^2]\simeq N^2+f^2(n\pi/kD)^2.
\]

\sourceart[0.54]{ChapmanRizzoli4.pdf}{17}{190bp 265bp 160bp 350bp}

Note that all waves have $\sigma>f$ and that all internal modes have $\sigma<N$. These two
limits are also points of vanishing $\vec c_g$ which is easily seen since $\partial\sigma/\partial k\to0$ there.

It is useful to examine the kinematics of the internal modes. We have
\[
 w=w_0e^{-i\sigma t+ikx}\sin[kR(z+D)].
\]
From continuity ($u_x+w_z=0$)
\[
 u=iRw_0e^{-i\sigma t+ikx}\cos[kR(z+D)].
\]

\sourcepagebreak

% -----------------------------------------------------------------------------
% Source printed page 81 / physical page 18
% -----------------------------------------------------------------------------
Consider the lowest mode $n=1$. Then from the dispersion relation $k_1\simeq\pi/RD$ and
\[
 w=w_0\cos(k_1x-\sigma t)\sin[\pi(z+D)/D]\ ;\qquad
 u=-Rw_0\sin(k_1x-\sigma t)\cos[\pi(z+D)/D]
\]
after taking the real part. The vertical structure looks like

\sourceart[0.70]{ChapmanRizzoli4.pdf}{18}{115bp 430bp 100bp 95bp}

Thus, we see that the particle motions under the crest and trough of a travelling wave
consists of a series of convergences and divergences giving a system of vertical cells.

\sourceart[0.78]{ChapmanRizzoli4.pdf}{18}{85bp 185bp 75bp 380bp}

One common consequence of this pattern is the formation of surface slicks or bands of
smooth, unrippled surface water, the bands being aligned parallel to the internal wave
crests. The scenario is as follows. A very thin organic film (one or two molecules thick)

\sourcepagebreak

% -----------------------------------------------------------------------------
% Source printed page 82 / physical page 19
% -----------------------------------------------------------------------------
typically covers the water surface. The periodic convergences and divergences of the
horizontal surface current due to the internal waves produce periodic contractions and
expansions of the surface film. This leads to an increase in the amount of film over the
convergences. The effect of the film in general is to reduce the surface tension, thereby
decreasing the tendency for short surface and capillary waves to form as the wind blows
over the water. Thus, the region over the convergences tends to have less ripples almost
to the point of elimination. These are the surface slicks. Such surface slicks have often
been used to infer the presence of internal waves, especially internal solitary waves.

\subsection{Evanescent modes}

We can reexamine the cases considered above but assuming that the wave decays in
the $x$ direction. For case A, we have
\[
 w=e^{-i\sigma t+kx}\omega(z)
\]
\[
 (\sigma^2-f^2)\omega_z+gk^2\omega=0\qquad\text{at }z=0
\]
\[
 \omega_{zz}+k^2R_1^2\omega=0
\]
\[
 \omega=0\qquad\text{at }z=-D.
\]
Solutions are
\[
 w=e^{-i\sigma t+kx}\sin[kR_1(z+D)]
\]
\[
 \frac{R_1(\sigma^2-f^2)}{gk}=-\tan kR_1D.
\]

\sourcepagebreak

% -----------------------------------------------------------------------------
% Source printed page 83 / physical page 20
% -----------------------------------------------------------------------------
\sourceart[0.77]{ChapmanRizzoli4.pdf}{20}{85bp 420bp 75bp 80bp}

There is an infinite set of evanescent modes although case A1 has two more than case
A. Putting on the rigid lid reduces A1 to A, i.e. $\sin(kR_1D)=0$.

For case B, we have
\[
 w=e^{-i\sigma t+kx}\omega(z)
\]
\[
 (\sigma^2-f^2)\omega_z+gk^2\omega=0\qquad\text{at }z=0
\]
\[
 \omega_{zz}-k^2R^2\omega=0
\]
\[
 \omega=0\qquad\text{at }z=-D.
\]
Solutions are
\[
 w=e^{-i\sigma t+kx}\sinh[kR(z+D)]
\]
\[
 \frac{R(\sigma^2-f^2)}{gk}=-\tanh kRD.
\]

\sourcepagebreak

% -----------------------------------------------------------------------------
% Source printed page 84 / physical page 21
% -----------------------------------------------------------------------------
\sourceart[0.74]{ChapmanRizzoli4.pdf}{21}{95bp 470bp 75bp 65bp}

We could have arrived at these same evanescent modes by setting $k=-ik'$ in the
previous section.

A summary of internal wave properties for all of the various ranges of
parameters can be found in Gill (1982, p. 261).

\section{Generation at a horizontal boundary}

We have discussed the properties of freely propagating internal gravity waves, but we
have not discussed how these waves might be generated in the ocean or the
atmosphere. One way that internal waves are generated is by a mean horizontal flow
passing over some topographic feature which forces the flow to move up and down
slightly. The following example illustrates this mechanism.

For simplicity, we neglect the effects of rotation and consider two-dimensional
flow over a sinusoidally varying horizontal boundary at $z=0$. The amplitude of the
variations is assumed small so that the dynamics can be linearized. Of course, an
arbitrarily shaped boundary could be used by first Fourier decomposing it, then
solving for the flow over each individual component, and then summing the results.

\setcounter{page}{85}

% -----------------------------------------------------------------------------
% Source printed page 85 / physical page 22
% -----------------------------------------------------------------------------
The mean flow has magnitude $U$ in the $x$ direction.

\sourceart[0.82]{ChapmanRizzoli4.pdf}{22}{70bp 500bp 55bp 80bp}

The topography has the form $h=h_0\sin(kx)$ with amplitude $h_0$. Moving along with
the mean flow, the topography has the form
\[
 h=h_0\sin[k(x+Ut)]
\]
from which we see that the frequency of the resulting motions will be
\[
 \sigma=-Uk.
\]
The topography introduces a vertical velocity because the particles near the boundary
must follow, to some extent, the undulations of the boundary. So,
\[
 w=U\frac{\partial h}{\partial x}=w_0e^{-i\sigma t+ikx}
 \qquad\text{on}\qquad z=0.
\]
This says that $w_0=Ukh_0$. The field equation for the region above the boundary is
\[
 w_{zz}-\frac{N^2-\sigma^2}{\sigma^2}w_{xx}=0.
\]
A solution is
\[
 w=w_0e^{-i\sigma t+ikx+imz}
\]
where
\[
 m^2=k^2(N^2-\sigma^2)/\sigma^2=(N/U)^2-k^2
\]

\sourcepagebreak

% -----------------------------------------------------------------------------
% Source printed page 86 / physical page 23
% -----------------------------------------------------------------------------
after substituting for $\sigma^2$. This solution satisfies the boundary condition at $z=0$ and
represents waves with phases propagating downward (because $\sigma<0$) which
corresponds to energy propagating upward to $z\to\infty$. Thus, the radiation condition at
$z\to\infty$ is satisfied, i.e. no energy enters the system from external sources. We now
examine two cases.

Suppose $\sigma^2>N^2$ which means $k>N/U$. This corresponds to short
wavelengths or undulations on the boundary. In this case $m^2<0$ so that $m$ must be
imaginary. The solution becomes
\[
 w=w_0e^{-i\sigma t+ikx-mz}\ ;\qquad
 m^2=k^2-(N/U)^2
\]
where the sign of $m$ is chosen to ensure that the solution remains finite. Recall that
\[
 \rho_0w_{zt}=p_{xx}
\]
from the momentum and continuity equations. This produces
\[
 \rho_0i\sigma mw=-k^2p
\]
from which
\[
 p=\frac{-\rho_0i\sigma m}{k^2}w
 =\frac{-\rho_0i\sigma m}{k^2}w_0e^{-i\sigma t+ikx-mz}.
\]
We see that $w$ and $p$ are out of phase by $\pi/2$. Therefore, the vertical energy flux,
$\overline{wp}=0$, is identically equal to zero, i.e. there is no vertical energy flux. This makes
sense because the solution decays exponentially in the vertical, so the waves cannot
transport any energy away from the boundary. Instead, the oscillations are trapped at
the boundary. If the wavelength is very small ($kU\gg N$), then stratification has little
effect and the flow is essentially irrotational.

Suppose $\sigma^2<N^2$ which means $k<N/U$. This corresponds to longer
wavelengths or undulations on the boundary. Now $m^2>0$, so $m$ is real and
$m^2=(N/U)^2-k^2$. The solution is
\[
 w=w_0e^{-i\sigma t+ikx+imz}.
\]

\sourcepagebreak

% -----------------------------------------------------------------------------
% Source printed page 87 / physical page 24
% -----------------------------------------------------------------------------
The form of this solution says that energy is continually being transported toward
$z\to\infty$. We have $\rho_0\sigma mw=-k^2p$, so
\[
 p=\frac{-\rho_0\sigma m}{k^2}w
 =\frac{-\rho_0\sigma m}{k^2}w_0e^{-i\sigma t+ikx+imz}.
\]
Now $w$ and $p$ are in phase, so $\overline{wp}\ne0$ and there is a net upward flux of energy. This
produces a drag on the mean flow because the energy must come from the mean flow.
The drag per unit surface area is the rate at which horizontal momentum is transferred
vertically
\[
 \tau=-\rho_0\overline{uw}
 =\frac{\overline{wp}}{U}
 =\frac{-\rho_0\sigma m w_0^2}{2k^2}\frac1U
 =\frac12\rho_0kh_0^2U^2
 \left(\frac{N^2}{U^2}-k^2\right)^{1/2}.
\]
The cutoff wavenumber which separates the two cases, $k_c=N/U$, corresponds
to the wavelength $2\pi/k_c$ which is the horizontal distance traveled by a particle in one
buoyancy period. This says that if the particle encounters multiple crests in the
topography during one buoyancy period, then the fluid will be forced to oscillate at
such a high frequency (greater than $N$) that no free waves can exist and no net drag
will be produced. If the particle stays within a single undulation, then the flow
adjustments will be slow enough so that free waves will be radiated away producing a
drag on the mean flow.

\section{Reflection from a solid boundary}

Here we consider the reflection of an internal gravity wave from a solid boundary which
is at some angle to the horizontal. To start, consider the two-dimensional solution
$e^{-i\sigma t+ikx+imz}$ which satisfies
\[
 w_{zz}-R^2w_{xx}=0
\]

\sourcepagebreak

% -----------------------------------------------------------------------------
% Source printed page 88 / physical page 25
% -----------------------------------------------------------------------------
where $R^2=(N^2-\sigma^2)/(\sigma^2-f^2)$ and $m=\pm Rk$. Lines of constant phase are those for
which
\[
 -\sigma t+kx\pm Rkz=\text{constant}.
\]
That is
\[
 x\pm Rz=(\sigma/k)t+\text{constant}.
\]

\sourceart[0.78]{ChapmanRizzoli4.pdf}{25}{85bp 390bp 70bp 150bp}

Phases propagate at right angles to $x\pm Rz=\text{constant}$. We see that energy flows
along $x\pm Rz=\text{constant}$ because
\[
 z_x=\pm R^{-1}
 =\pm\left(\frac{\sigma^2-f^2}{N^2-\sigma^2}\right)^{1/2}
 =\pm\left(\frac{(N^2-f^2)\cos^2\phi}
                 {(N^2-f^2)\sin^2\phi}\right)^{1/2}
 =\pm\cot\phi.
\]
These lines are the characteristics of the \emph{hyperbolic} $w$ equation, i.e.
\[
 w=F(x+Rz)+G(x-Rz)
\]
is the solution.

Now consider reflection from a solid plane wall passing through the origin;
$z=ax$. Remember that, for energy incident along $x+Rz=0$, the incident
wavenumber is along the normal to that line.

\sourcepagebreak

% -----------------------------------------------------------------------------
% Source printed page 89 / physical page 26
% -----------------------------------------------------------------------------
\sourceart[0.80]{ChapmanRizzoli4.pdf}{26}{75bp 435bp 55bp 100bp}

Energy exits along $x-Rz=0$, so the reflected wavenumber is normal to that line.
The frequency of the wave is determined solely by the angle to the vertical, and it
cannot change upon reflection. Thus the incident and reflected waves must make equal
angles with the rotation vector or the vertical (gravity) rather than with the normal to
the surface. Therefore, the reflection is not specular.

Consider the details of the situation. The incident wave is $\vec k_i$ and the reflected
wave is $\vec k_r$.

\sourceart[0.58]{ChapmanRizzoli4.pdf}{26}{170bp 210bp 150bp 345bp}

The projection of incident and reflected wavenumbers along $z=ax$ must be equal.
\[
 |\vec k_i|\cos(\tan^{-1}R-\tan^{-1}a)
 =|\vec k_r|\cos(\tan^{-1}R+\tan^{-1}a).
\]
We can evaluate these by geometry and the law of cosines
[$\cos\theta=(a^2+b^2-c^2)/2ab$].

\sourcepagebreak

% -----------------------------------------------------------------------------
% Source printed page 90 / physical page 27
% -----------------------------------------------------------------------------
\sourceart[0.68]{ChapmanRizzoli4.pdf}{27}{130bp 590bp 110bp 70bp}

After some algebra, the result is
\[
 |\vec k_r|=|\vec k_i|\left(\frac{1+aR}{1-aR}\right)
\]
\[
 m_r=\pm m_i\left(\frac{1+aR}{1-aR}\right)
 \left(\frac{R+a}{R-a}\right)
\]
where the signs have been taken from the sketch. This gives the new wavenumbers in
terms of the old ones. Because waves of a given frequency $\sigma$ can only go in the two
directions $\pm\tan^{-1}R$, reflection occurs not in the normal to the reflecting surface, but
rather in the direction of the stability gradient, i.e. in the $z$ direction, or in the
rotation vector. We have

\sourceart[0.72]{ChapmanRizzoli4.pdf}{27}{95bp 205bp 75bp 500bp}

The normal component of velocity must vanish at the wall, so
\[
 |\vec v_i|\sin(\tan^{-1}R^{-1}+\tan^{-1}a)
 =|\vec v_r|\sin(\tan^{-1}R^{-1}-\tan^{-1}a).
\]
Again we can use geometry to obtain
\[
 |\vec v_r|=|\vec v_i|\left(\frac{1+aR}{1-aR}\right).
\]

\sourcepagebreak

% -----------------------------------------------------------------------------
% Source printed page 91 / physical page 28
% -----------------------------------------------------------------------------
Notice that if $aR=1$, the reflected velocity is very large. What does this mean?
It means that the bottom coincides with the outgoing characteristic; $z=ax$ is the
bottom and the outgoing characteristic has the same slope. So, as $aR\to1$,
$|\vec v_r|$ becomes large and $\vec k_r$ becomes large, so that the reflected wave is very short. The present analysis
fails because we have neglected the effects of viscosity which would reduce the velocity
to zero at the boundary (no-slip condition), thereby avoiding the infinite $|\vec v_r|$.

We can visualize the reflection from various slopes, always requiring equal
projection of $|\vec v_i|$ and $|\vec v_r|$ on the normal to the surface.

\sourceart[0.80]{ChapmanRizzoli4.pdf}{28}{70bp 180bp 55bp 340bp}

\sourcepagebreak

% -----------------------------------------------------------------------------
% Source printed page 92 / physical page 29
% -----------------------------------------------------------------------------
\section{Variable buoyancy frequency}

We have restricted discussion to the case where the buoyancy frequency is constant,
i.e. constant vertical density gradient $\partial\rho/\partial z$. However, the more realistic situation is
when the buoyancy frequency varies with depth. Typical profiles of density and $N^2(z)$
in the ocean are

\sourceart[0.78]{ChapmanRizzoli4.pdf}{29}{85bp 350bp 75bp 130bp}

In most of the ocean $N^2>f^2$ although there are not many reliable values for the
deepest parts of the ocean. With this sort of profile, we have
\[
 R^2(z)=\frac{N^2(z)-\sigma^2}{\sigma^2-f^2}
\]
and $R^2$ may be greater than or less than zero. We can examine the changes in the
solution which result from this vertical dependence by looking for a wave solution like
\[
 w=e^{-i\sigma t+ikx}\omega(z).
\]
The waveguide internal wave problem becomes
\[
 (\sigma^2-f^2)\omega_z-gk^2\omega=0\qquad\text{at }z=0
\]
\[
 \omega_{zz}+k^2R^2(z)\omega=0
\]

\sourcepagebreak

% -----------------------------------------------------------------------------
% Source printed page 93 / physical page 30
% -----------------------------------------------------------------------------
\[
 \omega=0\qquad\text{at }z=-D.
\]
If $R^2(z)>0$, then the solution is trigonometric (travelling wave) because
$\omega_{zz}/\omega<0$. If $R^2(z)<0$, then the solution is exponential (evanescent mode) because
$\omega_{zz}/\omega>0$.

Suppose the $R^2(z)$ profile looks like

\sourceart[0.46]{ChapmanRizzoli4.pdf}{30}{195bp 390bp 175bp 160bp}

The system of equations is almost a Sturm--Liouville problem. We won't go into the
details of Sturm--Liouville theory because it can be found in many textbooks (and you
should be familiar with it). We can summarize some relevant properties which will be
useful for understanding the present problem. The usual Sturm--Liouville problem is
\[
 (p\psi_z)_z+(q+\lambda r)\psi=0
\]
\[
 a_{1,2}\psi_z+b_{1,2}\psi=0\qquad\text{at }z=0,-D
\]
with $p,r>0$; $q<0$ or $q>0$. The parameters $p,q,r,a,b$ are all real. There is a singly
infinite, denumerable set of eigenvalues
\[
 \lambda_1\le\lambda_2\le\lambda_3\le\cdots\to\infty.
\]
The eigenfunctions are orthogonal and can be orthonormalized as
\[
 \int_{-D}^{0}\psi_i\psi_jr(z)\,dz=\delta_{ij}
\]
where $\delta_{ij}$ is the Kronecker delta.

\sourcepagebreak

% -----------------------------------------------------------------------------
% Source printed page 94 / physical page 31
% -----------------------------------------------------------------------------
In terms of the Sturm--Liouville notation, we have
\[
 \lambda=k^2,\qquad r=R^2(z),\qquad p=1,\qquad q=0.
\]
Our present problem differs from the standard Sturm--Liouville problem in two ways.
First, the eigenvalue $k^2$ appears in the boundary condition at $z=0$. Second, $R^2(z)$
may change sign in $-D<z<0$. Let's look at each of these differences. When $R^2(z)$
changes sign in the definition interval, the generalized Sturm--Liouville theory
demonstrates the existence of an infinite, denumerable set of eigenvalues which are real
and have no lower bound to the sequence
\[
 -\infty<\cdots<(k_2^e)^2<(k_1^e)^2<0<(k_1^w)^2<(k_2^w)^2<\cdots<+\infty
\]
where the superscript $e$ represents an evanescent mode, and the superscript $w$
represents a travelling wave. The situation corresponds to

\sourceart[0.72]{ChapmanRizzoli4.pdf}{31}{105bp 305bp 90bp 240bp}

Thus, in this case, both evanescent and travelling modes are present simultaneously.

The effect of the appearance of $k^2$ in the surface boundary condition manifests
itself as
\[
 \int_{-D}^{0}\omega_i\omega_jR^2(z)\,dz
 =-\omega_i(0)\omega_j(0)g/(\sigma^2-f^2).
\]
This says that the eigenfunctions $\omega_i$ are not orthogonal unless the surface boundary is
the rigid lid, i.e. $\omega_i(0)=0$.

\sourcepagebreak

% -----------------------------------------------------------------------------
% Source printed page 95 / physical page 32
% -----------------------------------------------------------------------------
We could have solved the problem using the WKB approach in the case of
slowly varying $N^2(z)$, but we don't have time for that here. We can summarize the
situation for $R^2(z)$ as follows. If $R^2(z)$ does not change sign in the water column, then
all of the results found for the $N^2=\text{constant}$ case apply, with the only changes being
in the details of the dispersion relation and in the vertical dependence of the velocity
field. If $R^2(z)$ changes sign in the water column, then there is an infinite number of
evanescent modes and travelling waves present simultaneously. If $\sigma^2>f^2$, there are
also two travelling surface waves; if $\sigma^2<f^2$, there are two evanescent surface modes.
